\chapter{Conclusion}\label{chap:7}

Six pairs solved Tangram together in shared overhead VR and related physical and first-person conditions. This thesis followed how they judged shapes, expressed ideas through manipulation and responded to existing constructions.

The first finding concerns the geometric problems encountered during construction. The researcher's records identified near-complete shape problems in three of six first-person VR pairs and none of the six shared-overhead VR pairs. G06's candle showed why an accepted ending deserved closer examination: all seven pieces had been incorporated, yet the overall outline still differed from the target. This case connects the presentation of planar geometry to the pair's practical judgment of its result.

The second and third findings explain how constructions developed. Participants made local ideas visible through placements that their partners could extend or change. Work moved between separate parts, assistance and shared reconstruction. Partners also retained or continued partial arrangements that their original builders later revised. Self-revision could help the task progress or remove a correct part, while concrete alternatives could prompt direct partner intervention.

The project developed from eight proposed combinations and a separate construction task into Tangram across four conditions, with participants organising their collaboration freely. It delivered a shared overhead prototype and a six-pair corpus. Formal rounds used a fixed viewpoint. Target presentation changed between G01's initial recording and the revised sessions for G02--G06. These development decisions and the fixed condition order define the comparison presented here.

These processes support an inspectable overview, usable manipulation and brief assistance. Further evaluation can test viewpoint switching and access to earlier configurations with balanced tasks and stable controls.

